\clearpage
\begin{thebibliography}{99}
\bibitem{RiisTree}
S. Riis, \emph{A complexity gap for tree resolution},
Computational Complexity \textbf{10}, 179--209, 2001.
\href{https://doi.org/10.1007/s00037-001-8194-y}{doi:10.1007/s00037-001-8194-y}.

\bibitem{DantchevRiis}
S. Dantchev and S. Riis, \emph{On relativisation and complexity gap for
resolution-based proof systems}, in Computer Science Logic (CSL 2003),
Lecture Notes in Computer Science \textbf{2803}, 142--154, 2003.
\href{https://doi.org/10.1007/978-3-540-45220-1_14}{doi:10.1007/978-3-540-45220-1\_14}.

\bibitem{DMS}
S. Dantchev, B. Martin and S. Szeider, \emph{Parameterized proof complexity},
in Proceedings of the 48th Annual IEEE Symposium on Foundations of Computer Science
(FOCS 2007), 150--160, 2007.
\href{https://doi.org/10.1109/FOCS.2007.53}{doi:10.1109/FOCS.2007.53}.

\bibitem{DMRank}
S. Dantchev and B. Martin, \emph{Rank complexity gap for Lov\'asz--Schrijver
and Sherali--Adams proof systems}, Computational Complexity \textbf{22},
191--213, 2013.
\href{https://doi.org/10.1007/s00037-012-0049-1}{doi:10.1007/s00037-012-0049-1}.

\bibitem{QBFGap}
O. Beyersdorff, J. Clymo, S. Dantchev and B. Martin, \emph{The Riis complexity
gap for QBF resolution}, Journal on Satisfiability, Boolean Modeling and
Computation \textbf{15}(1), 9--25, 2024.
\href{https://doi.org/10.3233/SAT-231505}{doi:10.3233/SAT-231505}.

\bibitem{Riis}
S. Riis, \emph{On the asymptotic Nullstellensatz and Polynomial Calculus proof complexity}, in
Proceedings of the 23rd Annual IEEE Symposium on Logic in Computer Science
(LICS 2008), pp.~272--283, 2008.
\href{https://doi.org/10.1109/LICS.2008.30}{doi:10.1109/LICS.2008.30}.

\bibitem{RiisLong}
S. Riis, \emph{On the asymptotic Nullstellensatz and Polynomial Calculus proof complexity}, extended preprint,
Electronic Colloquium on Computational Complexity, Report TR09-014, 2009.
\url{https://eccc.weizmann.ac.il/report/2009/014/}.

\bibitem{Krajicek}
J. Kraj\'i\v cek, \emph{On the degree of ideal membership proofs from uniform
families of polynomials over a finite field}, Illinois Journal of
Mathematics \textbf{45}(1), 41--73, 2001.
\href{https://doi.org/10.1215/ijm/1258138254}{doi:10.1215/ijm/1258138254}.

\bibitem{Harman}
N. Harman, \emph{Stability and periodicity in the modular representation
theory of symmetric groups}, preprint, arXiv:1509.06414.
\url{https://arxiv.org/abs/1509.06414}.

\bibitem{HKPSYZ}
C. Hone, F. Klein, B. Pauwels, A. Sherman, O. Yacobi and V. L. Zhang,
\emph{Semisimplifying categorical Heisenberg actions and periodic equivalences}, preprint,
arXiv:2509.07377.
\url{https://arxiv.org/abs/2509.07377}.

\bibitem{BDM}
C. Bowman, S. Doty and S. Martin, \emph{Integral Schur--Weyl duality for
partition algebras}, Algebraic Combinatorics \textbf{5}(2), 371--399, 2022.
\href{https://doi.org/10.5802/alco.214}{doi:10.5802/alco.214}.
\bibitem{KrajEuler}
J. Kraj\'i\v cek, \emph{Uniform families of polynomial equations over a
finite field and structures admitting an Euler characteristic of definable
sets}, Proceedings of the London Mathematical Society (3)
\textbf{81}(2), 257--284, 2000.
\href{https://doi.org/10.1112/S0024611500012557}{doi:10.1112/S0024611500012557}.

\bibitem{KrajScanlon}
J. Kraj\'i\v cek and T. Scanlon, \emph{Combinatorics with definable sets:
Euler characteristics and Grothendieck rings}, Bulletin of Symbolic Logic
\textbf{6}(3), 311--330, 2000.
\href{https://doi.org/10.2307/421058}{doi:10.2307/421058}.

\bibitem{CEF}
T. Church, J. S. Ellenberg and B. Farb, \emph{FI-modules and stability
for representations of symmetric groups}, Duke Mathematical Journal
\textbf{164}(9), 1833--1910, 2015.
\href{https://doi.org/10.1215/00127094-3120274}{doi:10.1215/00127094-3120274}.
\bibitem{LangAlgebra}
S. Lang, \emph{Algebra}, revised third edition,
Graduate Texts in Mathematics \textbf{211}, Springer, 2002.
\href{https://doi.org/10.1007/978-1-4613-0041-0}{doi:10.1007/978-1-4613-0041-0}.

\bibitem{AlperinLocal}
J. L. Alperin, \emph{Local Representation Theory:
Modular Representations as an Introduction to the Local
Representation Theory of Finite Groups},
Cambridge Studies in Advanced Mathematics \textbf{11},
Cambridge University Press, 1986.
\href{https://doi.org/10.1017/CBO9780511623592}{doi:10.1017/CBO9780511623592}.

\bibitem{Cutland}
N. Cutland, \emph{Computability: An Introduction to Recursive
Function Theory}, Cambridge University Press, 1980.
\href{https://doi.org/10.1017/CBO9781139171496}{doi:10.1017/CBO9781139171496}.

\bibitem{StacksExact}
The Stacks Project Authors,
Lemma 12.7.2 (Tag 010N) in \emph{Additive functors}, and
Remark 35.4.4 (Tag 08WK) in \emph{Descent for universally injective morphisms}.
\url{https://stacks.math.columbia.edu/tag/010N};
\url{https://stacks.math.columbia.edu/tag/08WK}.

\bibitem{ParisWilkie}
J. Paris and A. Wilkie, \emph{Counting problems in bounded arithmetic},
in C. A. Di Prisco (ed.), \emph{Methods in Mathematical Logic},
Lecture Notes in Mathematics \textbf{1130}, 317--340, Springer, 1985.
\href{https://doi.org/10.1007/BFb0075316}{doi:10.1007/BFb0075316}.


\bibitem{NSFoundations}
P. Beame, R. Impagliazzo, J. Kraj\'i\v cek, T. Pitassi and P. Pudl\'ak,
\emph{Lower bounds on Hilbert's Nullstellensatz and propositional proofs},
Proceedings of the London Mathematical Society (3)
\textbf{73}(1), 1--26, 1996.
\href{https://doi.org/10.1112/plms/s3-73.1.1}{doi:10.1112/plms/s3-73.1.1}.
Preliminary version in FOCS~1994, 794--806.

\bibitem{PCFoundations}
M. Clegg, J. Edmonds and R. Impagliazzo,
\emph{Using the Groebner basis algorithm to find proofs of unsatisfiability},
in Proceedings of the 28th Annual ACM Symposium on Theory of Computing
(STOC~1996), 174--183, 1996.
\href{https://doi.org/10.1145/237814.237860}{doi:10.1145/237814.237860}.

\bibitem{RiisSitharamTautologies}
S. Riis and M. Sitharam,
\emph{Generating hard tautologies using predicate logic and the symmetric group},
Logic Journal of the IGPL \textbf{8}(6), 787--795, 2000.
\href{https://doi.org/10.1093/jigpal/8.6.787}{doi:10.1093/jigpal/8.6.787}.

\bibitem{RiisSitharamModules}
S. Riis and M. Sitharam,
\emph{Uniformly generated submodules of permutation modules:
Over fields of characteristic 0},
Journal of Pure and Applied Algebra \textbf{160}(2--3), 285--318, 2001.
\href{https://doi.org/10.1016/S0022-4049(00)00069-4}{doi:10.1016/S0022-4049(00)00069-4}.

\bibitem{Ajtai}
M. Ajtai, \emph{Symmetric systems of linear equations modulo p},
Electronic Colloquium on Computational Complexity, Report TR94-015, 1994.
\url{https://eccc.weizmann.ac.il/report/1994/015/}.

\bibitem{BussDesign}
S. R. Buss, \emph{Lower bounds on Nullstellensatz proofs via designs},
in P. Beame and S. R. Buss (eds.), \emph{Proof Complexity and Feasible
Arithmetics}, DIMACS Series in Discrete Mathematics and Theoretical
Computer Science \textbf{39}, 59--71, American Mathematical Society, 1998.
\url{https://mathweb.ucsd.edu/~sbuss/ResearchWeb/designs/}.

\bibitem{BeameRiis}
P. Beame and S. Riis, \emph{More on the relative strength of counting
principles}, in P. Beame and S. R. Buss (eds.), \emph{Proof Complexity
and Feasible Arithmetics}, DIMACS Series in Discrete Mathematics and
Theoretical Computer Science \textbf{39}, 13--35,
American Mathematical Society, 1998.
\url{https://homes.cs.washington.edu/~beame/papers/riis.pdf}.

\bibitem{PTT}
F. Part, N. Thapen and I. Tzameret, \emph{First-order reasoning and
efficient semi-algebraic proofs}, Annals of Pure and Applied Logic
\textbf{176}(1), 103496, 2025.
\href{https://doi.org/10.1016/j.apal.2024.103496}{doi:10.1016/j.apal.2024.103496}.

\bibitem{RiisSitharam1997}
S. Riis and M. Sitharam, \emph{Non-constant degree lower bounds imply
linear degree lower bounds}, Electronic Colloquium on Computational
Complexity, Report TR97-048, 1997.
\url{https://eccc.weizmann.ac.il/report/1997/048/}.
The proof qualification is discussed in the introduction.

\bibitem{RiisSitharamTranslationPreprint}
S. Riis and M. Sitharam, \emph{Generating hard tautologies using predicate
logic and the symmetric group}, BRICS Report RS-98-19, 1998.
Earlier version of~\cite{RiisSitharamTautologies}; Section~II,
Theorem~1 gives the Skolemised translation and model--solution correspondence.
\url{https://www.brics.dk/RS/98/19/BRICS-RS-98-19.pdf}.

\bibitem{RiisSupporting}
S.~Riis.
\emph{Degree profiles of uniform polynomial consequences: Linear lower
bounds for Nullstellensatz and Polynomial Calculus---supporting materials}.
2026. Paper, tutorial, Lean formalisation, verification records and
reproduction instructions. GitHub snapshot, 4 November 2026,
commit \texttt{072d949}.
\href{https://github.com/SR123/degree-profiles-linear-bounds/archive/072d949368e3ad7f44fc3b34b80d8f43fa92dcd1.zip}{Archive (ZIP)};
\href{https://github.com/SR123/degree-profiles-linear-bounds/blob/072d949368e3ad7f44fc3b34b80d8f43fa92dcd1/paper/paper.pdf}{archived paper (PDF)};
\href{https://github.com/SR123/degree-profiles-linear-bounds/tree/072d949368e3ad7f44fc3b34b80d8f43fa92dcd1/lean}{Lean sources and verification}.

\bibitem{RiisThesis}
S. Riis, \emph{Independence in Bounded Arithmetic},
Ph.D. thesis, University of Oxford, 1993.

\bibitem{RiisFinitisation}
S. Riis, \emph{Finitisation in Bounded Arithmetic},
BRICS Report Series RS-94-23, University of Aarhus, 1994.
\url{https://doi.org/10.7146/brics.v1i23.21644}.

\bibitem{KrajBounded}
J. Kraj\'{i}\v{c}ek, \emph{Bounded Arithmetic, Propositional Logic,
and Complexity Theory}, Encyclopedia of Mathematics and its Applications,
vol.~60, Cambridge University Press, 1995, Chapter~9.

\bibitem{RazborovPC}
A. A. Razborov, \emph{Lower bounds for the polynomial calculus},
Computational Complexity \textbf{7}(4), 291--324, 1998.
\url{https://doi.org/10.1007/s000370050013}.

\end{thebibliography}
